\documentclass[11pt,a4paper]{article}
\usepackage[utf8]{inputenc}

% Math
\usepackage{amsmath, amssymb, amsthm}

% Tables and lists
\usepackage{booktabs}
\usepackage{array}
\usepackage{enumitem}

% Design
\usepackage{xcolor}
\usepackage[most]{tcolorbox}
\usepackage{geometry}
\geometry{margin=2.5cm}

% ============================================================
% Theme: change these three colours to re-theme the document.
% ============================================================
\definecolor{accent}{HTML}{8B2B40}
\definecolor{accentDark}{HTML}{5D3A1A}
\definecolor{accentLight}{HTML}{F9F4E8}

% ============================================================
% Content boxes: the toolbar's "Boxes" menu inserts these styles.
% ============================================================
\tcbset{
    definitionBox/.style={
        enhanced, breakable, colback=accentLight, colframe=accent,
        fonttitle=\bfseries, coltitle=white, colbacktitle=accent,
        boxrule=1pt, arc=4pt,
        left=10pt, right=10pt, top=5pt, bottom=5pt,
    },
    noteBox/.style={
        enhanced, breakable, colback=white, colframe=accent,
        fonttitle=\bfseries, coltitle=accent, colbacktitle=white,
        boxrule=0pt, leftrule=4pt, sharp corners,
        left=10pt, right=10pt,
    },
    exampleBox/.style={
        enhanced, breakable, colback=accentLight!60!white, colframe=accentDark,
        fonttitle=\bfseries, coltitle=white,
        boxrule=0.5pt, arc=3pt,
        left=10pt, right=10pt,
    },
    textBox/.style={
        enhanced, breakable, colback=white, colframe=accentDark,
        boxrule=0pt, leftrule=2pt, sharp corners,
        left=10pt, right=10pt,
    },
    genericBox/.style={
        enhanced, breakable, colback=accentLight!60!white, colframe=accentDark,
        boxrule=0.5pt, arc=3pt,
        left=10pt, right=10pt,
    },
}

% Theorem-like environments
\theoremstyle{plain}
\newtheorem{theorem}{Theorem}[section]
\theoremstyle{remark}
\newtheorem*{remark}{Remark}

% Coloured list markers
\setlist[itemize,1]{label={\color{accent}\textbullet}}
\setlist[enumerate,1]{label={\color{accent}\bfseries\arabic*.}}

% A simple title block
\newcommand{\notesTitle}[2]{%
    \noindent{\color{accent}\rule{\linewidth}{3pt}}\par
    \vspace{1.2em}
    \begin{center}
        {\color{accentDark}\LARGE\bfseries #1}\par
        \vspace{0.5em}
        {\color{accent}\large #2}\par
        \vspace{0.4em}
        {\small \today}
    \end{center}
    \vspace{0.6em}
    \noindent{\color{accent}\rule{\linewidth}{1pt}}\par
    \vspace{1.5em}
}

% hyperref goes last: it patches many other packages.
\usepackage{xurl}
\usepackage[
    colorlinks=true,
    urlcolor=accent,
    linkcolor=accent,
    breaklinks=true,
    pdfpagemode=UseNone
]{hyperref}

\begin{document}

\notesTitle{{{title}}}{Linear Algebra --- Eigenvalues and Eigenvectors}

\section{Motivation}

Many linear transformations have special directions that they only \emph{stretch}, never rotate. Finding those directions turns a hard matrix problem into a handful of scalar ones. Press \textbf{Ctrl+Enter} to compile, and scroll either pane: the other follows, anchored on the section headings.

\begin{tcolorbox}[definitionBox, title=Definition 1.1]
Let $A$ be an $n \times n$ matrix. A nonzero vector $\mathbf{v}$ is an \emph{eigenvector} of $A$ if
\[ A\mathbf{v} = \lambda \mathbf{v} \]
for some scalar $\lambda$, called the corresponding \emph{eigenvalue}.
\end{tcolorbox}

\section{Finding eigenvalues}

Rewriting the definition as $(A - \lambda I)\mathbf{v} = \mathbf{0}$ shows that a nonzero solution exists exactly when $A - \lambda I$ is singular.

\begin{theorem}
$\lambda$ is an eigenvalue of $A$ if and only if
\[ \det(A - \lambda I) = 0. \]
\end{theorem}

\begin{tcolorbox}[noteBox, title=Note]
The polynomial $p(\lambda) = \det(A - \lambda I)$ is the \emph{characteristic polynomial}. It has degree $n$, so $A$ has at most $n$ distinct eigenvalues.
\end{tcolorbox}

\subsection{A worked example}

\begin{tcolorbox}[exampleBox, title=Example 2.1]
Let
\[ A = \begin{pmatrix} 2 & 1 \\ 1 & 2 \end{pmatrix}. \]
Then
\[ \det(A - \lambda I) = (2-\lambda)^2 - 1 = (\lambda - 1)(\lambda - 3), \]
so the eigenvalues are $\lambda_1 = 1$ and $\lambda_2 = 3$, with eigenvectors $(1,-1)^{T}$ and $(1,1)^{T}$.
\end{tcolorbox}

\section{Summary}

\begin{table}[h]
    \centering
    \begin{tabular}{@{}lll@{}}
        \toprule
        \textbf{Concept} & \textbf{Equation} & \textbf{Meaning} \\
        \midrule
        Eigenpair            & $A\mathbf{v} = \lambda\mathbf{v}$  & Direction only scaled \\
        Characteristic eq.   & $\det(A - \lambda I) = 0$          & Finds every $\lambda$ \\
        Trace                & $\operatorname{tr} A = \sum \lambda_i$ & Sum of eigenvalues \\
        Determinant          & $\det A = \prod \lambda_i$         & Product of eigenvalues \\
        \bottomrule
    \end{tabular}
\end{table}

\begin{remark}
Symmetric real matrices, like the one in Example 2.1, always have real eigenvalues and orthogonal eigenvectors.
\end{remark}

\section{Exercises}

\begin{enumerate}
    \item Find the eigenvalues of $\begin{pmatrix} 4 & 0 \\ 0 & -1 \end{pmatrix}$.
    \item Which of the following is the characteristic polynomial of $\begin{pmatrix} 1 & 2 \\ 0 & 3 \end{pmatrix}$?
    \begin{enumerate}[label=\alph*),leftmargin=1.5em,topsep=0.3em,itemsep=0.15em]
        \item $\lambda^2 - 4\lambda + 3$
        \item $\lambda^2 + 4\lambda + 3$
        \item $\lambda^2 - 3$
    \end{enumerate}
\end{enumerate}

\vspace{1em}
\noindent Further reading: \href{https://en.wikipedia.org/wiki/Eigenvalues_and_eigenvectors}{Eigenvalues and eigenvectors}.

\end{document}
